\documentclass[11pt,a4paper,]{article}
\usepackage{lmodern}
\usepackage{amssymb,amsmath}
\usepackage{ifxetex,ifluatex}
\usepackage{fixltx2e} % provides \textsubscript
\ifnum 0\ifxetex 1\fi\ifluatex 1\fi=0 % if pdftex
  \usepackage[T1]{fontenc}
  \usepackage[utf8]{inputenc}
\else % if luatex or xelatex
  \ifxetex
    \usepackage{mathspec}
    \usepackage{xltxtra,xunicode}
  \else
    \usepackage{fontspec}
  \fi
  \defaultfontfeatures{Mapping=tex-text,Scale=MatchLowercase}
  \newcommand{\euro}{€}
\fi
% use upquote if available, for straight quotes in verbatim environments
\IfFileExists{upquote.sty}{\usepackage{upquote}}{}
% use microtype if available
\IfFileExists{microtype.sty}{%
\usepackage{microtype}
\UseMicrotypeSet[protrusion]{basicmath} % disable protrusion for tt fonts
}{}
\usepackage[margin=2cm]{geometry}
\ifxetex
  \usepackage[setpagesize=false, % page size defined by xetex
              unicode=false, % unicode breaks when used with xetex
              xetex]{hyperref}
\else
  \usepackage[unicode=true]{hyperref}
\fi
\hypersetup{breaklinks=true,
            bookmarks=true,
            pdfauthor={},
            pdftitle={Untitled},
            colorlinks=true,
            citecolor=blue,
            urlcolor=blue,
            linkcolor=magenta,
            pdfborder={0 0 0}}
\urlstyle{same}  % don't use monospace font for urls
\setlength{\parindent}{0pt}
\setlength{\parskip}{6pt plus 2pt minus 1pt}
\setlength{\emergencystretch}{3em}  % prevent overfull lines
\setcounter{secnumdepth}{0}

%%% Use protect on footnotes to avoid problems with footnotes in titles
\let\rmarkdownfootnote\footnote%
\def\footnote{\protect\rmarkdownfootnote}


\title{
\noindent
\large\textbf{Untitled} \hfill \textbf{FirstName LastName} \\
\normalsize Module: ModuleNumber \hfill Team members: student1, student2 \\
Module Lead: Prof.~Smith \hfill Lab Date: XX/XX/XX \\
Lab Lead: John Smith \hfill Due Date: XX/XX/XX
}
\date{\vspace{-5ex}}


\begin{document}

\maketitle





\section{About}\label{about}

This template is similar to the
\texttt{03:\ My\ Template\ (HTML\ edit)}, but instead applies a custom
LaTeX template. This is used to build the final output.

\subsection{Alteration to the
template}\label{alteration-to-the-template}

You will notice that there are many additional YAML options than usual.
At line 153, the template has been altered to include a lot of extra
fields, which are enclosed by the dollar signs:

\begin{verbatim}
\title{
\noindent
\large\textbf{$title$} \hfill \textbf{$first-name$ $last-name$} \\
\normalsize Module: $module$ \hfill Team members: $for(teammates)$$teammates$$sep$, $endfor$ \\
Module Lead: $module-lead$ \hfill Lab Date: $lab-date$ \\
Lab Lead: $TA$ \hfill Due Date: $due-date$
}
\date{\vspace{-5ex}}
\end{verbatim}

\section{Example feature.}\label{example-feature.}

This template also includes a bibliography file, containing a single
reference for R Markdown (Allaire et al. 2018). The bibliography will be
added to the end of the output document.

\subsection{Customisation}\label{customisation}

If you want to create your own LaTeX template, you will have to be at
least partly familiar with LaTeX. But this example highlights how the
title section can be customised to include additional fields beyond the
typical \texttt{title}, \texttt{date} and \texttt{authors} etc.

\section*{References}\label{references}
\addcontentsline{toc}{section}{References}

\phantomsection\label{refs}
\begin{CSLReferences}{1}{0}
\bibitem[\citeproctext]{ref-R-rmarkdown}
Allaire, JJ, Yihui Xie, Jonathan McPherson, Javier Luraschi, Kevin
Ushey, Aron Atkins, Hadley Wickham, Joe Cheng, and Winston Chang. 2018.
\emph{Rmarkdown: Dynamic Documents for r}.
\url{https://CRAN.R-project.org/package=rmarkdown}.

\end{CSLReferences}

\end{document}
